\section{Confirmatory results}\label{sec:confirm}

Everything in this section uses the \NTest\ test scenarios only (\NTestPairs\ pairs). The hypotheses,
thresholds and decision rule are the preregistered ones; the estimator (resamples, seed, quality filter,
verdict rule) was fixed afterwards in the confirmatory script (\cref{sec:design}). The results below depend on
that estimator, so this section is labelled confirmatory but not preregistered. The preregistration said the cost ratio is computed ``among pairs where quality is
non-inferior.'' That sentence has two readings, so both are reported:
\begin{itemize}
  \item \textbf{Pair reading (co-primary):} keep a pair only if its own turn-pass fraction is not worse than
    its anchor's by more than \NIMarginPts\ points. In this design that keeps pairs where the arm did at
    least as well. Because this filter conditions on an outcome the arm affects, it is not reported alone.
  \item \textbf{Arm reading (co-primary, unfiltered):} keep every cost-valid pair, and require non-inferiority of
    the arm as a whole.
\end{itemize}

\begin{table}[htbp]
\centering\small
\caption{Primary endpoint: geometric-mean cost ratio, arm over plain host, test split, with 95\,\%
scenario-cluster bootstrap intervals. ``Pairs kept'' counts pairs that pass the pair-level quality filter
out of all cost-valid test pairs.}
\label{tab:primary}
\begin{tabular*}{\textwidth}{@{\extracolsep{\fill}}l l r r r r r@{}}
\toprule
 & & \multicolumn{3}{c}{Pair reading (co-primary)} & \multicolumn{2}{c}{Arm reading (co-primary, unfiltered)} \\
\cmidrule(lr){3-5}\cmidrule(l){6-7}
Host & Arm & pairs kept & ratio & 95\,\% CI & ratio & 95\,\% CI \\
\midrule
Fable 5.1 & sticky & \CfFableStickyPairN/\CfFableStickyNValid & \CfFableStickyPair & \CfFableStickyPairCI & \CfFableStickyArm & \CfFableStickyArmCI \\
Fable 5.1 & shipped & \CfFableShippedPairN/\CfFableShippedNValid & \CfFableShippedPair & \CfFableShippedPairCI & \CfFableShippedArm & \CfFableShippedArmCI \\
Fable 5.1 & sonnet & \CfFableSonnetPairN/\CfFableSonnetNValid & \CfFableSonnetPair & \CfFableSonnetPairCI & \CfFableSonnetArm & \CfFableSonnetArmCI \\
\addlinespace[3pt]
Opus 5.5 & sticky & \CfOpusStickyPairN/\CfOpusStickyNValid & \CfOpusStickyPair & \CfOpusStickyPairCI & \CfOpusStickyArm & \CfOpusStickyArmCI \\
Opus 5.5 & shipped & \CfOpusShippedPairN/\CfOpusShippedNValid & \CfOpusShippedPair & \CfOpusShippedPairCI & \CfOpusShippedArm & \CfOpusShippedArmCI \\
Opus 5.5 & sonnet & \CfOpusSonnetPairN/\CfOpusSonnetNValid & \CfOpusSonnetPair & \CfOpusSonnetPairCI & \CfOpusSonnetArm & \CfOpusSonnetArmCI \\
\bottomrule
\end{tabular*}

\end{table}
\howtoreadtable{A ratio below 1 is a saving: \CfFableStickyPair\ means the routed arm paid about
\CfFableStickyCents\ cents for every dollar the plain host paid. If the whole interval is below 1, the
cost difference is not explained by chance; if it is above 1, neither is the extra cost. The pair-level filter drops sessions
that were cheaper but did worse, so for sticky and shipped it raises the ratio by \FilterLiftLow\ to
\FilterLiftHigh; on Opus the arm reading gives \CfOpusStickyArmExtra--\CfOpusSonnetArmExtra\,\% more
than the plain host.}

\begin{figure}[htbp]
\centering
\pgfplotslegendfromname{forestlegend}\par\smallskip
% Confirmatory cost ratios (test split) with 95% scenario-cluster bootstrap CIs. Log x axis.
\begin{tikzpicture}
\begin{axis}[benchmark, xmode=log, log ticks with fixed point,
  xtick={0.5,0.6,0.7,0.8,0.9,1,1.2,1.4,1.6}, xmin=0.45, xmax=1.75,
  scale only axis, width=10.5cm, xmajorgrids, y tick label style={font=\small}, height=5.2cm, ymin=-0.7, ymax=6.7,
  ytick={0,1,2,4,5,6}, yticklabels from table={generated/data/forest.dat}{label},
  xlabel={Cost ratio, routed arm / plain host (log scale; below 1 saves money)},
  legend to name=forestlegend, legend columns=2, legend cell align=left,
  legend style={/tikz/every even column/.append style={column sep=10pt}}]
\draw[black!70, line width=0.8pt] (axis cs:1,-0.7) -- (axis cs:1,6.7);
\draw[okgray, dashed, line width=0.8pt] (axis cs:\HTwoLower,-0.7) -- (axis cs:\HTwoLower,2.6);
\addplot[only marks, mark=*, mark size=2.6pt, color=okblue,
  error bars/.cd, x dir=both, x explicit, error bar style={line width=1pt}]
  table[x=pair, y=y, x error minus=pm, x error plus=pp]{generated/data/forest.dat};
\addlegendentry{pair reading (co-primary)}
\addplot[only marks, mark=diamond, mark size=3pt, color=okorange, mark options={line width=1pt},
  error bars/.cd, x dir=both, x explicit, error bar style={line width=0.6pt, densely dashed}]
  table[x=arm, y expr=\thisrow{y}-0.3, x error minus=am, x error plus=ap]{generated/data/forest.dat};
\addlegendentry{arm reading (co-primary, unfiltered)}
\end{axis}
\end{tikzpicture}

\smallskip
\begin{tikzpicture}
\begin{axis}[benchmark, xmode=log, log ticks with fixed point,
  xtick={0.5,0.6,0.7,0.8,0.9,1,1.2,1.4,1.6}, xmin=0.45, xmax=1.75,
  scale only axis, width=10.5cm, xmajorgrids, y tick label style={font=\small}, height=1.6cm, ymin=-0.7, ymax=1.7,
  ytick={0,1}, yticklabels from table={generated/data/forest-h3.dat}{label},
  xlabel={Cost ratio, sticky / shipped on the same scenario (H3; log scale)}]
\draw[black!70, line width=0.8pt] (axis cs:1,-0.7) -- (axis cs:1,1.7);
\draw[okgray, dashed, line width=0.8pt] (axis cs:\HThreeUpper,-0.7) -- (axis cs:\HThreeUpper,1.7);
\addplot[only marks, mark=*, mark size=2.6pt, color=okblue,
  error bars/.cd, x dir=both, x explicit, error bar style={line width=1pt}]
  table[x=pair, y=y, x error minus=pm, x error plus=pp]{generated/data/forest-h3.dat};
\addplot[only marks, mark=diamond, mark size=3pt, color=okorange, mark options={line width=1pt},
  error bars/.cd, x dir=both, x explicit, error bar style={line width=0.6pt, densely dashed}]
  table[x=arm, y expr=\thisrow{y}-0.3, x error minus=am, x error plus=ap]{generated/data/forest-h3.dat};
\end{axis}
\end{tikzpicture}

\caption{Confirmatory cost ratios on the test split. Top: each routed arm against the plain host. Bottom: the
sticky arm against the shipped arm on the same scenarios (H3). Whiskers are 95\,\% intervals. The solid line
marks 1 (no difference); the dashed lines mark the preregistered thresholds for H2 (\HTwoLower) and H3
(\HThreeUpper).}
\label{fig:forest}
\howtoread{Each row is one comparison; the dot is the best estimate and the whisker the 95\,\% interval.
The horizontal axis is logarithmic, so ``half the cost'' and ``double the cost'' are the same distance from
1. All three Fable rows lie well left of 1 (savings); all three Opus rows lie right of the dashed line
(no saving). In the bottom panel both hosts lie left of 1: the sticky policy cost less than the shipped one, by a
small margin on Fable.}
\end{figure}

\subsection{The hypotheses}

\begin{table}[htbp]
\centering\small
\caption{The preregistered hypotheses, test split, pair reading. The arm reading gives the same verdicts.}
\label{tab:hyp}
\begin{tabular*}{\textwidth}{@{\extracolsep{\fill}}>{\raggedright\arraybackslash}p{0.33\textwidth} l l r r l@{}}
\toprule
Hypothesis (test split, pair reading) & Host & Arm & ratio & 95\,\% CI & verdict \\
\midrule
H1: routing is cheaper (upper end $<$ 1) & Fable & sticky & \CfFableStickyPair & \CfFableStickyPairCI & \vconf \\
H1: routing is cheaper (upper end $<$ 1) & Fable & shipped & \CfFableShippedPair & \CfFableShippedPairCI & \vconf \\
H2: no saving (lower end $>$ \HTwoLower) & Opus & sticky & \CfOpusStickyPair & \CfOpusStickyPairCI & \vconf \\
 & Opus & shipped & \CfOpusShippedPair & \CfOpusShippedPairCI & \vconf \\
 & Opus & sonnet & \CfOpusSonnetPair & \CfOpusSonnetPairCI & \vconf \\
H3: sticky/shipped (upper end $\le$ \HThreeUpper) & Fable & sticky/shipped & \HThreeFable & \HThreeFableCI & \vconf \\
 & Opus & sticky/shipped & \HThreeOpus & \HThreeOpusCI & \vconf \\
\bottomrule
\end{tabular*}

\end{table}

\begin{itemize}
  \item \textbf{H1, confirmed.} On the Fable host, sticky cost \CfFableStickyPair\X\ (95\,\% CI
    \CfFableStickyPairCI) and shipped \CfFableShippedPair\X\ (\CfFableShippedPairCI) of the plain host.
  \item \textbf{H2, confirmed.} On the Opus host, sticky cost \CfOpusStickyPair\X\ (\CfOpusStickyPairCI),
    shipped \CfOpusShippedPair\X\ (\CfOpusShippedPairCI) and plain Sonnet \CfOpusSonnetPair\X\
    (\CfOpusSonnetPairCI): routing to Sonnet costs about \CfOpusStickyExtra--\CfOpusSonnetExtra\,\% more.
  \item \textbf{H3, confirmed.} The preregistered test was only that sticky costs no more than shipped
    (upper end at most \HThreeUpper). Sticky cost \HThreeFable\X\ of shipped on Fable (\HThreeFableCI) and
    \HThreeOpus\X\ on Opus (\HThreeOpusCI), on \HThreeFableN\ and \HThreeOpusN\ matched scenario
    repetitions. On Fable the upper end, \HThreeFableUpper, is close to 1.
\end{itemize}

\subsection{Quality}

\begin{table}[htbp]
\centering\small
\caption{Quality: mean difference in turn-pass fraction, arm minus plain host, test split, all quality-valid
pairs. Non-inferior when the lower end of the interval is above \NIMargin.}
\label{tab:quality}
\begin{tabular*}{\textwidth}{@{\extracolsep{\fill}}l l r r r l@{}}
\toprule
Host & Arm & pairs & mean $\Delta$ turn-pass & 95\,\% CI & non-inferior \\
\midrule
Fable 5.1 & sticky & 46 & \QFableSticky & \QFableStickyCI & \vconf \\
Fable 5.1 & shipped & 46 & \QFableShipped & \QFableShippedCI & \vconf \\
Fable 5.1 & sonnet & 46 & \QFableSonnet & \QFableSonnetCI & \vconf \\
Opus 5.5 & sticky & 46 & \QOpusSticky & \QOpusStickyCI & \textcolor{okorange}{\textsf{not confirmed}} \\
Opus 5.5 & shipped & 46 & \QOpusShipped & \QOpusShippedCI & \vconf \\
Opus 5.5 & sonnet & 46 & \QOpusSonnet & \QOpusSonnetCI & \vconf \\
\bottomrule
\end{tabular*}

\end{table}

Every Fable savings claim passes the quality test. (A \term{savings claim} here is any arm and host whose
own cost interval shows a saving; only those cells are tested for quality.) The closest call is Fable sticky,
whose lower bound is \QFableStickyLow, inside the \NIMargin\ margin by only \QFableStickyMarginGap.
Across all \QAllScen\ scenarios (exploratory) its mean difference is \QAllFableSticky\ with a 95\,\%
interval of \QAllFableStickyCI: there the margin is not shown, and on the \NTrain\ training scenarios it is
\RvQTrainFableSticky\ (\RvQTrainFableStickyCI). So quality passed on the test split only. The quality test itself
is unfiltered: it uses every quality-valid test pair, so the \QFableStickyMarginGap\ margin does not depend on the
pair-level cost filter. That margin is close to the Monte-Carlo noise of the bootstrap: re-drawing the resamples
with a different random stream moves the lower bound by \RvQNoise\ (to \RvQAltLow). Opus sticky did not show
non-inferior quality (\QOpusSticky, 95\,\% CI \QOpusStickyCI; ``not shown'', which is not the same as ``worse'').
The preregistration requires quality for each savings claim, and no savings claim is made on Opus, so the quality
hypothesis is confirmed under that scope; if quality were required of every routed cell, \RvNQualCells\ of
\RvNRoutedCells\ cells would pass and \RvNHypAllQuality\ of 4 hypotheses would be confirmed.

\subsection{The savings model}\label{sec:modelcheck}

The preregistered deliverable is a model that predicts the saving of a new session from what is known before
it starts: arm, host, task type, number of turns, long pauses and the size of the first prompt. It was fitted
on the training scenarios only and then asked to predict the \MCPairs\ held-out test pairs.

\begin{figure}[htbp]
\centering
\pgfplotslegendfromname{modellegend}\par\smallskip
% Prediction model check on the test split: predicted total saving (with 90% PI) vs observed, per arm and host.
\begin{tikzpicture}
\begin{axis}[benchmark, scale only axis, width=\ModelAxisW, height=\ModelAxisH,
  xmin=-70, xmax=150, ymin=-70, ymax=150,
  xtick={-50,0,50,100,150}, ytick={-50,0,50,100,150},
  xlabel={Predicted total saving on the test pairs (\$; whisker: 90\,\% interval)},
  ylabel={Observed total saving (\$)}, xmajorgrids, ymajorgrids, clip=false,
  legend to name=modellegend, legend columns=3, legend cell align=left,
  legend style={/tikz/every even column/.append style={column sep=10pt}}]
\addplot[okgray, dashed, line width=0.8pt, domain=-70:150, samples=2] {x};
\addlegendentry{prediction = observation}
\addplot[only marks, mark=*, mark size=2.6pt, color=okblue,
  error bars/.cd, x dir=both, x explicit, error bar style={okblue!60, line width=0.7pt}]
  table[x=pred, y=obs, x error minus=minus, x error plus=plus]{generated/data/model-check-fable.dat};
\addlegendentry{Fable host}
\addplot[only marks, mark=square*, mark size=2.6pt, color=okorange,
  error bars/.cd, x dir=both, x explicit, error bar style={okorange!60, line width=0.7pt}]
  table[x=pred, y=obs, x error minus=minus, x error plus=plus]{generated/data/model-check-opus.dat};
\addlegendentry{Opus host}
% Generated by build_assets.py: labels for panel model-check. Do not edit.
\node[scatterlabel, anchor=north] at ($(axis cs:93.5700,85.0100)+(0.000pt,-4.700pt)$) {shipped, Fable};
\node[scatterlabel, anchor=north west] at ($(axis cs:-32.6400,-31.4800)+(3.323pt,-3.323pt)$) {shipped, Opus};
\node[scatterlabel, anchor=east] at ($(axis cs:69.1800,94.9500)+(-4.700pt,0.000pt)$) {sonnet, Fable};
\node[scatterlabel, anchor=north] at ($(axis cs:-42.5100,-46.8000)+(0.000pt,-4.700pt)$) {sonnet, Opus};
\node[scatterlabel, anchor=west] at ($(axis cs:103.1300,99.8200)+(4.700pt,0.000pt)$) {sticky, Fable};
\node[scatterlabel, anchor=west] at ($(axis cs:-28.6900,-21.1500)+(4.700pt,0.000pt)$) {sticky, Opus};

\end{axis}
\end{tikzpicture}

\caption{The savings model on the test split: predicted total saving (with its 90\,\% interval) against the
observed total, for each arm and host.}
\label{fig:model}
\howtoread{A point on the dashed diagonal is a perfect prediction. Every whisker crosses the diagonal, so
every observed total lies inside its predicted 90\,\% interval. Points above zero are savings (Fable);
points below zero are extra costs (Opus).}
\end{figure}

The model passed both preregistered checks. Its 90\,\% intervals covered \MCCovLog\ of test pairs on the log
ratio and \MCCovUSD\ on the dollar saving, both inside [\CovLow, \CovHigh]. The observed test total,
\$\MCObserved, lay inside the predicted interval of \$\MCPILow\ to \$\MCPIHigh\ (point prediction
\$\MCPredicted). The intervals are conservative: they covered \MCCovLog\ of pairs against a nominal
\MCNominal. The pooled total nets Fable savings (\$\MCFableTotal) against Opus losses (\$\MCOpusTotal), so
it says little about either host on its own. The design document's stricter criteria (coverage
\DesignCovLow--\DesignCovHigh, calibration slope \DesignSlopeLow--\DesignSlopeHigh) are also met (slope
\MCSlope). \Cref{tab:modelcheck} in the appendix gives the check for each arm and host. The verdict is
``\MCVerdict''; its limits are discussed in \cref{sec:modellimits}.

\subsection{Decision}

Applying the preregistered rule: on the Fable host, \textbf{\DecisionFable}; on the Opus host,
\textbf{\DecisionOpus}. Every verdict was re-derived under both non-inferiority readings and under
\NSeeds\ other bootstrap seeds (\SeedLo--\SeedHi); \SeedsChanged\ verdicts changed. The other
clarifications of the preregistration (\cref{sec:deviations}) were not re-run separately.
